\chapter[Can Reason Redesign Society?]{Can Human Beings Redesign Society through Reason Alone?}
\section{A Storm in a White Porcelain Cup}
Pick up a small white porcelain coffee cup, thin-walled and gold-rimmed, made in eighteenth-century Paris or London.

Its contents have travelled farther than the cup. Coffee originated in Ethiopia's highlands; Arabs first cultivated it in Yemen; Ottoman ports carried it to Istanbul; Venetian merchants brought it into Europe. By the late seventeenth century, cafés filled the streets of Paris, London and Amsterdam. A penny admitted you to a London coffeehouse, with newspapers to read and debates to hear. Hence their nickname: ``penny universities'', admitting students for one coin.

The nickname hides a bold idea: learning and public affairs no longer belong exclusively to churches, courts and universities. A tailor, sailor or apothecary can sit down for a penny, offer opinions on royal policy, war or a new book, and expect a serious hearing.

The cup thus becomes a weapon, aimed not at someone's head but at an entire order. Why organise society by birth, churches and inherited rules? Why not reason it out from first principles like a geometry problem, then rebuild according to the plan?

Those staking everything on this question became known as Enlightenment thinkers; their period became the Enlightenment. Our question is their wager: can people redesign society through reason alone?

\section{A Paris Café around 1750}
Come into a scene.

Evening in Paris, at the Procope on the Rue de l'Ancienne-Comédie. Open for nearly seventy years, it is among the city's best-known cafés. Open the door and smells greet you: fresh coffee, pipe tobacco, candle smoke and customers' damp winter wool. Small round marble tables stand close together; candles cast shadows on the ceiling. Racks hold the day's newspapers and new pamphlets for anyone to read and replace.

In a corner, a thin middle-aged man writes rapidly beside stacked manuscript pages: Denis Diderot, rushing entries for the Encyclopédie he edits. Its almost arrogant ambition is to collect all existing human knowledge, from theology to handicrafts, arrange it alphabetically and sell it. Eventually reaching twenty-eight volumes, it was banned by the French government and placed on the Church's Index, becoming more famous with each prohibition. Knowledge ceased to be locked in monasteries and palaces. It could sit in bookshops for anyone able to pay. That alone was revolutionary.

Regulars know the window table as Voltaire's favourite. Stories credit him with dozens of coffees daily, probably exaggerated, but the seat is real. He is likely in Ferney near the Swiss border now, while his books slip into Paris in hidden cart compartments.

Farther inside, two customers argue red-faced while others laughingly intervene. Jean-Jacques Rousseau appeared on evenings like this too, before falling out with almost everyone.

One person goes unnoticed: the young waiter threading between tables with a tray. While guests discuss reason belonging to everyone, nobody asks him. He cannot read; newspapers are merely patterned paper. Remember him. Later we shall repeatedly return to people standing outside the café's circle of light.

\section{An Old Carpet or a New Blueprint?}
Set out the conflict before answering.

To Enlightenment thinkers, society resembled an inherited carpet, centuries old, faded and patched with hierarchy, privilege and ``as it has always been''. Nobles enjoyed tax exemptions while peasants and townspeople shouldered almost everything. Books required a censor's approval. An artisan's son's possibilities were largely assigned at birth. Patches upon patches had been sewn not according to reason, but because they had always been there.

The startling proposal: pull away the carpet and weave a new one from a blueprint. Drawn with what? Reason. Bring everything---law, taxation, education, government and religious organisation---before reason and demand its justification. What cannot answer has no right to remain.

Two cracks immediately open. The rest of our chapter inhabits them.

First: is reason enough? Blueprint advocates ask why geometry can calculate a bridge's load but not a good law. Newton explained planets through a few laws; why not society through a few principles? Opponents answer that society is not a bridge. A miscalculated bridge does not hate you; people do. Every old patch may preserve a forgotten lesson. Before removing it, do you know why it exists?

Second, sharper: even if a blueprint is possible, who draws it? Voltaire hoped for an enlightened monarch, a king who read philosophy. Rousseau said only the people themselves. But who counts as the people? The illiterate waiter? Women? Enslaved plantation workers across the Atlantic? This second crack eventually made the whole Enlightenment creak.

\section{One Table, Many Positions}
First avoid a misunderstanding. ``The Enlightenment'' sounds like a unified movement with a programme. It was more like interconnected debating contests across Europe and North America, with centres in Edinburgh, Königsberg, Milan, Naples and Philadelphia as well as Paris. Participants quarrelled with one another as fiercely as with the old order. They agreed not on everything, but that everything should be open to argument.

Over decades of debate, they accumulated shared possessions still present in your life:

First, \textbf{reason}: every claim requires reasons; ancestral practice or a book's authority is not enough.

Second, \textbf{toleration}: differences of faith or opinion should not bring execution, imprisonment or confiscation. This agreement was bought with blood. Over the preceding two centuries, repeated wars between Catholics and Protestants killed or wounded millions. Toleration was not politeness, but a ceasefire.

Third, \textbf{natural rights}: certain rights come with being human, not from kings or parliaments. In the Two Treatises of Government, Locke argued, broadly, that government rests on the governed's consent and exists to protect life, liberty and property. Systematic violation gives the people the right to replace it.

Fourth, \textbf{public discussion}: arguments belong not only in studies and palaces, but openly in coffeehouses, periodicals and salons, available for everyone's objections.

Fifth, \textbf{progress}: humanity need not decline generation after generation or circle endlessly. Knowledge accumulates, institutions improve, and descendants may live better than ancestors. This was startling when most people assumed antiquity lay closer to perfection.

Their disagreements mattered as much as their common ground.

Voltaire chiefly fought religious persecution. In Toulouse in 1762, the Protestant merchant Jean Calas was accused of murdering a son who supposedly intended to become Catholic. Despite highly questionable evidence, he was executed by being torn apart. Voltaire spent years investigating and campaigning, published the Treatise on Tolerance and secured exoneration in 1765. Those are facts. His political design looks less attractive: distrusting direct popular rule, he hoped philosophy could persuade enlightened monarchs, corresponding with and visiting Prussian and Russian rulers.

Montesquieu offered another blueprint in The Spirit of the Laws: liberty cannot depend on rulers' kindness, but on powers checking one another. Separate legislative, executive and judicial authority so nobody absorbs all. He credited observations of English politics for the inspiration.

Rousseau disagreed with both. Sovereignty belonged to the people collectively; every scheme delegating it would deteriorate. He and Voltaire broke publicly and exchanged sarcasm to the end. Ironically, textbooks now print all three together on one page.

Another counterexample guards against error. You may equate Enlightenment with worship of omnipotent reason. Yet David Hume, a central Scottish Enlightenment figure in Edinburgh, wrote in A Treatise of Human Nature: ``Reason is, and ought only to be the slave of the passions.'' Reason calculates routes; feelings and desires determine destinations. At the movement's centre sat someone drawing reason's limits. Treating an age as a solid block is almost always wrong.

Finally remember a woman at the meeting's edge. Madame Geoffrin hosted a celebrated Paris salon, regularly attended by Diderot and fellow Encyclopédie leaders. Scholars throughout Europe prized invitations. She could determine the subject and participants, yet could not place her name on the Encyclopédie she helped make possible. What separated the lit table from the unlit space? Hold that question; we shall return.

\section[Thinking Bridge: The State of Nature]{Thinking Bridge: A ``State of Nature'' Thought Experiment}
Enlightenment thinkers had a useful tool you can apply in class meetings or family disagreements: the \textbf{thought experiment}. Spend no money, move no bricks. Build an imagined scene, push the problem to its extreme and see what follows.

A favourite was the state of nature: imagine a nonexistent moment with no government, laws, police or schools. How would humans live?

Its cleverness lies in replacing an unobservable question---why society is as it is---with a tractable one: what would we lack without society? Those absences reveal what society supplies.

Three thinkers obtained three results from the same experiment.

Hobbes had lived through the English Civil War and seen collapsed order's horrors. Broadly, he reasoned that everyone in nature pursues security, pre-emptive attack pays, and mutual suspicion creates a war of all against all. People must surrender most freedoms for a government too strong to defy. The more terrible nature appears, the stronger government's justification.

Locke's version is milder. People naturally possess reason and moral sense, knowing not to kill or rob. Their problem is the lack of an impartial judge: everyone judges their own case and revenge becomes endless. Government supplies not security itself, but fair arbitration. An unfair judge breaks the contract and can be replaced.

Rousseau's version is most radical. Natural humans are free, self-sufficient and mutually unindebted. Where would war arise? Society creates inequality. Once someone encloses land, declares it theirs and finds believers, hierarchy, property and enslavement accumulate.

See the manoeuvre? \textbf{Each quietly places his conclusion inside his starting assumptions.} Naturally selfish, aggressive people yield strong government; reasonable people yield limited government; naturally good people yield an indictment of society. Thought experiments expose reasoning but cannot prove initial premises. Those require evidence and judgements about human nature, always open to dispute.

A descendant survives today. In A Theory of Justice, twentieth-century American philosopher John Rawls proposed a veil of ignorance. Before designing rules, imagine not knowing your future wealth, gender, health or abilities. Rules you would then accept qualify as fair. This repeats the old method: an imagined scene makes ``What if I were that person?'' open to reasoning.

Next time you assess a school rule, classroom rule or platform policy, do not start with liking or disliking. Imagine its absence. Then imagine occupying its worst-off position. Would you still agree?

\section{Two Passages from Two Centuries Ago}
Read a primary source. In 1784, a Prussian periodical invited answers to ``What is Enlightenment?'' Fifty-nine-year-old Kant, the Königsberg philosopher whose regular walks supposedly let neighbours set their clocks, wrote a short response opening:

\begin{quote}
``Enlightenment is humanity's emergence from its self-imposed immaturity. \ldots{} Have the courage to use your own understanding! That is the Enlightenment's motto.'' (Kant, ``An Answer to the Question: What Is Enlightenment?'', 1784.)
\end{quote}
Notice the seven-character Chinese phrase rendered ``self-imposed''. Kant's immaturity is not insufficient intelligence: humans naturally have understanding. It is laziness or fear in using it, preferring books, priests or officials to think for us. Some chains are handed over voluntarily. Enlightenment therefore begins as a problem of courage, not knowledge. He also distinguishes two uses of reason. As officer, taxpayer or employee, one obeys one's duties. Outside that role, as a person, one may publicly question any institution. He calls this latter activity reason's public use. Remember it for our discussion of the public sphere.

Now the first page of Rousseau's Social Contract, 1762:

\begin{quote}
``Man is born free, and everywhere he is in chains.'' (Rousseau, The Social Contract, Book I, chapter 1.)
\end{quote}
Its fame comes precisely from its inconsistency. If people are born free, whence the chains? If chains are everywhere, what does born free mean? Rousseau does not dissolve the contradiction. He makes it a lever under every existing institution: you claim authority over me---on what grounds?

Together the passages convey the age's temperature: confidence in humanity---you possess understanding, use it---and anger at the present---chains everywhere. Confidence plus anger fuelled Enlightenment.

\section{``All Are Born Equal''---What about Women?}
Now face the hardest question. Establish the facts first.

In 1789, the French Revolution began, and the National Constituent Assembly issued the Declaration of the Rights of Man and of the Citizen. Its first article reads:

\begin{quote}
``Men are born and remain free and equal in rights.'' (Declaration of the Rights of Man and of the Citizen, 1789, Article 1.)
\end{quote}
Thirteen years earlier, the North American Declaration of Independence offered the more resounding: ``We hold these truths to be self-evident, that all men are created equal.''

Around the same people, at the same time, French women were legally subordinate to fathers and husbands, without votes or independent property control. Atlantic slave trading was at its height. Slave-grown cane and coffee made Saint-Domingue France's most profitable colony. Jefferson, who wrote that all were created equal, held hundreds of enslaved people. Locke, arguing government existed to protect property, was deeply involved in colonial affairs and held shares in the slave-trading Royal African Company.

Separate our categories again. Declarations appeared, slavery persisted and women were excluded: evidence, indisputable. Why these coexisted is explanation, where arguments follow. Contemporary understandings differed: some sincerely saw no contradiction, some saw it and stayed silent, some shouted. Evaluation is your task; do not rush it.

Excluded people did not wait. They took the declarations' own language and asked: does this word ``human'' include me?

In 1791, Olympe de Gouges published the Declaration of the Rights of Woman and of the Female Citizen, deliberately echoing the earlier declaration article by article. Its opening says, broadly, that woman is born free and equal to man in rights. She later died by guillotine.

That same year, Saint-Domingue's enslaved people rose. They and their successors claimed freedom through the French Revolution's own language: liberty, equality and natural rights. After over a decade of war, Haiti declared independence in 1804, becoming the first republic founded by a slave uprising.

In 1792, Mary Wollstonecraft's A Vindication of the Rights of Woman argued that women's apparent lesser rationality arose not naturally but from denied education and training as pleasing ornaments. Give them equal schooling, then compare.

Compare three explanations for universal rights coexisting with universal exclusion.

\textbf{Explanation one: hypocrisy.} They never meant it. Fine declarations concealed clear calculations: property, colonies and plantation profits must remain untouched. Locke's shares and Jefferson's estate provide solid evidence. Yet this cannot explain excluded people's enthusiasm. If the declaration was only deception, why did de Gouges imitate it or Haitian rebels march under its language? A fraud would not become a weapon turned against its perpetrators.

\textbf{Explanation two: expansion.} Once uttered, universal language acquires its own life. Designers may intend a limited group, but cannot prevent outsiders reading literally and demanding fulfilment. Rights resemble seeds planted for apples that grow into a forest. Later abolition, women's rights and anticolonial movements repeatedly borrowed the language, supporting this view. But expansion was never automatic. Haiti paid in prolonged war, followed by French demands for huge compensation. French women gained voting rights only in 1944. Seeds do not spontaneously become forests; every growth ring required struggle. Crediting concepts alone erases fighters' blood and labour.

\textbf{Explanation three: exclusion was built in.} The universal person was originally modelled as an adult, property-owning white man. Exclusion was not faulty execution, but a hidden compartment in the blueprint. This explains its systematic, untroubled character: people were not overlooked, but never counted. Yet if exclusion were wholly intrinsic, excluded people should not have been able to use the language. History shows the reverse: it became their commonest weapon.

Each holds part of the truth. Hypocrisy directs attention to interests, expansion to language's power, and embedded exclusion to conceptual design. Keep all three. When reading any document beginning ``everyone'', ``all people'' or ``each person'', ask whose image shaped its human being.

\section[Reason's Critics]{When Reason Turns: Counter-Enlightenment and Romanticism}
The French Revolution first offered Enlightenment blueprints large-scale construction. Aristocratic privilege, guilds and Church estates were dismantled. Even measures and calendars seemed insufficiently rational and were redesigned. The metric system arose here, a genuine rational achievement still used worldwide. By 1793, however, construction changed character: invoking virtue and public safety, revolutionary authorities executed growing lists of enemies. Yesterday's declaration writers faced today's guillotine. This became the Terror.

Why declarations led to execution remains disputed, beyond our chapter's adjudication. File it under competing causal explanations. One fact is clear: before the Terror, in 1790, MP Edmund Burke published Reflections on the Revolution in France, predicting that blueprint revolution would slide into violence and military dictatorship. Subsequent events made him famous.

Now steelman the opposition: \textbf{make Burke's and the conservatives' case as fully as possible}. Reducing them to reactionary vested interests is unfair and wastes an argument deserving attention.

First, old institutions compress experience. Every carpet patch may record a forgotten lesson: a strange prohibition once answered an unknown epidemic, an elaborate ritual resolved a communal conflict. However elegant, a rational blueprint contains a few years of one mind's thought; custom contains generations of real trial and error. Before demolition, establish why each patch was sewn.

Second, complex systems do not accept blueprints. Society resembles an ecosystem, not a bridge. One change sends consequences through invisible channels. Since nobody foresees everything, reversible small improvements are safer than starting over.

Third, society is not a contract among contemporaries, but a partnership of the dead, living and unborn. Tradition contains predecessors' wisdom and sacrifices. Our generation has no right to sell it off: that steals from partners not yet born.

Strong arguments? Very. They explain post-revolutionary turmoil and imported good institutions failing to take root. But expose the fatal weakness: if all patches have histories, what happens to sufferers? A patch may preserve a lesson, or merely a nail driven by former rulers. ``Take it slowly'' tells plantation workers their children and grandchildren must remain slaves. Conservatism naturally favours those already adapted to current arrangements, its ineradicable blind spot. Before the 190s, women in no country possessed full political rights. If we always wait until the time is ripe, it may never ripen itself.

Burke's contemporary Herder objected from another direction. Enlightenment thinkers, he argued, spoke of humanity while imagining Paris salon-goers. They measured the world with one ruler, assigned humanity childhood, youth and adulthood, and announced themselves adults. For Herder, each people's language, songs and stories possessed a soul. Culture was a garden, not a ladder, each flower blooming differently. This bore good fruit in serious collection of folk songs, epics and languages, and poisonous fruit when later nationalism fermented difference into superiority.

Romanticism followed. Goethe's The Sorrows of Young Werther, published in 1774, described a sensitive young man's suicide over hopeless love. Europe was enthralled; young readers copied his blue coat and yellow waistcoat. Romantics restored neglected feeling, imagination, nature, mystery and nostalgia to the table's centre. Intriguingly, Rousseau became a precursor: designer of popular sovereignty in Enlightenment genealogies, singer of emotion and nature in Romantic ones. Labels come afterwards; living people resist neat compartments.

The Romantic question was not opposition to reason, but where emotion, belonging, beauty and sacredness should live in a designed society. Blueprints can draw fair taxes; can they draw the feeling of home? Enlightenment did not answer well, and neither have we.

\section[Further Challenge: The Public Sphere]{Further Challenge: What Did the Coffeehouse Become?}
Now introduce our weightiest theoretical tool. In 1962, Jürgen Habermas published The Structural Transformation of the Public Sphere, transforming our opening café into a sociological concept: the \textbf{public sphere}.

An everyday example helps. Discussing classroom matters at your dinner table is private. A headteacher announcing discipline at an assembly represents state power. Between them lies a space where ordinary people meet, read shared newspapers and debate public affairs: reasonable policies, good books, officials who should resign. That is the public sphere. Habermas argued that eighteenth-century coffeehouses, salons, reading societies and periodicals first made public opinion a political force. Previously power needed obedience; afterwards it began needing reasons.

This explains apparently scattered phenomena. Why did encyclopaedias and pamphlets provoke bans, burning and pursuit of booksellers? The old order understood the danger. Royal authority rested on questions not being asked; the public sphere let everyone ask. Why did Kant distinguish reason's public use? He recognised the new discussion space as Enlightenment's engine. One person privately questioning a king is a grumbler; ten thousand doing so in print and coffeehouses constitute an age.

The theory's most valuable feature, however, is its cracks, exposed by later researchers. Feminist historians revisited this golden age of equal discussion and found unequal admission. Geoffrin could host Europe's celebrated salon without signing the publications she enabled. Women, poor people, the illiterate and colonised people remained outside the coffeehouse door. Kant's public reason was an attractive emptiness for Procope's illiterate waiter: not because he lacked understanding, but because he could not enter that room.

Use the concept in two halves. One is a \textbf{measuring rod}: a healthy society needs an accessible, reason-governed discussion space between household and power. The other is a \textbf{revealing mirror}: whenever an age declares free speech for everyone, ask whether it truly means everyone. Who is outside? Whose voice is audible but not taken seriously?

A measuring rod and a revealing mirror: take them into your own era.

\section[Today's Designers]{Today: Blueprint Engineers Have Changed Professions}
The three-hundred-year-old wager never remained confined to Europe.

During nineteenth-century Japan's Meiji transformation, Fukuzawa Yukichi opened An Encouragement of Learning in 1872 with ``Heaven creates no person above another, nor any person below another'', restating equality in terms Japanese readers understood. He and colleagues formed the Meirokusha, publishing a journal and debating politics, giving Tokyo something of Paris salon life. In China around this period, Yan Fu translated Western works including Evolution and Ethics, introducing evolution, liberty and democracy into Chinese. Decades later, New Culture youth invited ``Mr De'', Democracy, and ``Mr Sai'', Science, to re-examine inherited rules. Every society meeting Enlightenment faced the same questions: remove the old carpet? How much? Who decides? Japan offered one answer, China several. None can claim the task finished.

Look around now. Those proposing rational social redesign have changed professions. Mostly they neither wear wigs nor inhabit cafés; they write code. Data designs urban traffic, algorithms rank news for 1.4 billion people, and scoring systems shape couriers', drivers' and merchants' fortunes. Behind them stands the old belief that society can be calculated and calculation can improve it. Technological optimism directly descends from progress, even sounding similar: a little more data and computing power will bring a better society.

Measure everyday life with the public sphere: is the internet a new coffeehouse?

Similarities are plain. Admission is almost free. Where an illiterate waiter could not enter Procope, a mountain-region secondary student can comment beneath the same post as a professor. Enlightenment thinkers scarcely dreamed of such equality.

Differences are more troubling. Coffeehouse participants jointly set discussion's pace. Online, algorithms decide visibility, judged by time spent rather than reason. Nearby tables could challenge café arguments; online echo chambers return one's own views. Kantian public reason requires exposing arguments to everyone's objections. Trending debates often assign camps before reasons appear. Tools improve without removing difficulties. ``Who may enter?'' becomes ``Who may be heard?''

Smaller echoes surround you. Who sets classroom rules? A teacher alone follows the enlightened-monarch model. A class vote raises another question: what if implementation harms a minority? Should cleaning duties follow lottery, rotation or consideration for distant commuters? Each is a miniature state-of-nature experiment, inviting Rousseau's question: you say this rule serves everyone---on what grounds?

\section[Your Turn: Pull Away the Carpet?]{Your Turn: Dare You Pull Away the Old Carpet?}
Return to the white porcelain cup.

Three centuries ago, café-goers believed conversation and reason could take society from ancestors' hands, dismantle it and redraw it. You use their inheritance daily: metric measures, public schools, no conviction without trial and the legitimacy of asking why you should obey. Their unfulfilled promises remain too: two centuries later, the declarations' everyone still admits new people.

Take the wager. Given one chance to redesign a small society around you---class, club or family---through reason alone, what would you do?

You have the blueprint case: public rules, equality, reasons for everything and accountable power. You have seen how ``always done this way'' shelters injustice.

You also have the patchwork case: old rules may preserve unknown lessons; complex systems resist perfect plans; belonging, ritual and indefinable affection elude calculation and may not be reassembled after demolition.

Which irrational old patch would you keep? Which inherited rule would you firmly remove? Choose and explain, with one requirement: you must be willing to show your reasons to someone in the worst-off position and read them aloud while meeting their eyes.

Three centuries ago, they urged courage in using your understanding. Perhaps you must finish their sentence: afterwards, have the courage to answer for the consequences.
